\section{Numerical demonstrations}
\label{sec:numerical}

We include two minimal exact-diagonalization checks that illustrate the
commuting result and the operator-algebra obstruction. These are not intended
as high-precision numerical studies; they are small finite-bath sanity checks
that reproduce the qualitative structure discussed above. The scripts are in
\texttt{simulations/src/main.py} with a convenience runner in
\texttt{simulations/run.ps1}. The default ``safe'' profile uses small truncations
and single-threaded linear algebra to reduce resource usage.

\paragraph{Commuting QND benchmark.}
We take a four-level system with diagonal
$H_Q=\mathrm{diag}(0,0.8,1.6,2.4)$ and $f=\mathrm{diag}(0,1,2,3)$, coupled to a
finite set of harmonic bath modes. For each inverse temperature $\beta$ we
compute $\rho_Q\propto\Tr_X e^{-\beta H_{\mathrm{tot}}}$ by exact
matrix exponentiation and compare it with the commuting prediction
$e^{-\beta H_Q}\exp[\tfrac12 C(\beta) f^2]$ where
$C(\beta)=\beta\sum_k c_k^2/\omega_k^2$ (for $m_k=1$). The script reports the
relative Frobenius error in \texttt{results/commuting\_check.csv}. Residual
error is dominated by bath truncation and decreases with larger oscillator
cutoffs.

\paragraph{Noncommuting qubit check.}
We take a qubit with $H_Q=0.6\,\sigma_z$ and
coupling $f=\sigma_x+\eta\,\sigma_z$ (with $\eta=0.3$ in the default profile),
again coupled to a single truncated oscillator. We compute
$H_{\mathrm{MF}}(\beta)=-(1/\beta)\log\rho_Q$ and decompose it in the Pauli basis.
For the commuting control $f=\sigma_z$ only the $\sigma_z$ component survives,
whereas the noncommuting case produces a finite transverse component (reported
in \texttt{results/qubit\_pauli\_coeffs.json}). This explicitly illustrates the
operator growth induced by noncommutativity.
